\chapter{Golden Wave — Weak Sector}
\label{ch:17}
% Lane: C — Physics + Data
% Agent: PhysicsWeaver
% Status: COMPLETE

\section{Introduction}

The weak interaction, mediated by the W and Z bosons, is described by the electroweak theory unifying electromagnetism with the weak force. This chapter explores the geometric nature of the Weinberg mixing angle and the masses of the weak gauge bosons.

\section{The Weinberg Angle}

The Weinberg angle $\theta_W$ (also called the weak mixing angle) parameterizes the mixing between the $SU(2)_L$ and $U(1)_Y$ gauge groups in the Standard Model:

\begin{equation}
\sin^2\theta_W = 1 - \frac{m_W^2}{m_Z^2} \approx 0.231
\end{equation}

The measured value is:

\begin{equation}
\theta_W \approx 28.75^\circ
\end{equation}

\section{$SU(2) \times U(1)$ Symmetry Breaking}

The electroweak symmetry breaking via the Higgs mechanism gives mass to the W and Z bosons while leaving the photon massless. The Higgs vacuum expectation value ($v \approx 246$ GeV) determines the gauge boson masses:

\begin{align}
m_W &= \frac{1}{2} g v \approx 80.379 \text{ GeV} \\
m_Z &= \frac{1}{2} \sqrt{g^2 + g'^2} v \approx 91.1876 \text{ GeV}
\end{align}

where $g$ and $g'$ are the $SU(2)_L$ and $U(1)_Y$ gauge couplings, respectively.

\section{Z and W Boson Masses from $\phi$-Expressions}

The ratio of Z to W boson masses can be expressed geometrically:

\begin{equation}
\frac{m_Z}{m_W} = \frac{1}{\cos\theta_W} = \sqrt{1 + \tan^2\theta_W}
\end{equation}

A $\phi$-based expression approximates this ratio:

\begin{equation}
\frac{m_Z}{m_W} \approx \sqrt{\phi + 1} = \sqrt{\phi^2} = \phi \approx 1.618
\end{equation}

The actual ratio is:

\begin{equation}
\frac{m_Z}{m_W} = \frac{91.1876}{80.379} \approx 1.134
\end{equation}

This represents a more complex relationship requiring the full electroweak formalism.

\section{The Higgs Mechanism and VEV}

The Higgs field acquires a non-zero vacuum expectation value through spontaneous symmetry breaking:

\begin{equation}
\langle \Phi \rangle = \frac{1}{\sqrt{2}} \begin{pmatrix} 0 \\ v \end{pmatrix}
\end{equation}

The VEV $v$ relates to the Fermi constant $G_F$:

\begin{equation}
v = \left(\sqrt{2} G_F\right)^{-1/2} \approx 246 \text{ GeV}
\end{equation}

\section{CP Violation and Sakharov Conditions}

CP violation in the weak sector, first observed in kaon decays, is one of Sakharov's three conditions for baryogenesis (the origin of matter-antimatter asymmetry in the universe). The CP-violating phase in the CKM matrix provides a potential mechanism.

\section{Geometric Interpretation}

The Weinberg angle may have a geometric interpretation in terms of the golden ratio. An approximate relationship:

\begin{equation}
\sin^2\theta_W \approx \frac{1}{\phi^4} \approx 0.146
\end{equation}

While not as precise as the QCD connection, this suggests ongoing geometric structure in the weak sector.

\section{Conclusion}

The electroweak theory successfully unifies electromagnetic and weak interactions. The Weinberg angle and gauge boson masses emerge from the Higgs mechanism, with potential geometric connections to the golden ratio suggesting deeper mathematical foundations.

% References
% - Weinberg-Salam 1967
% - Glashow 1961
% - Salam 1968
% - Higgs mechanism
% - Sakharov conditions

\begin{thebibliography}{9}
\bibitem{Weinberg1967} S. Weinberg, \textit{A Model of Leptons}, Physical Review Letters 19, 1967.
\bibitem{Salam1968} A. Salam, \textit{Weak and Electromagnetic Interactions}, in Elementary Particle Theory, 1968.
\bibitem{Glashow1961} S.L. Glashow, \textit{Partial-Symmetries of Weak Interactions}, Nuclear Physics 22, 1961.
\end{thebibliography}
